\uuid{HMrd}
\exo7id{592}
\titre{exo7 592}
\auteur{gourio}
\organisation{exo7}
\datecreate{2001-09-01}
\isIndication{false}
\isCorrection{false}
\chapitre{Suite}
\sousChapitre{Suites équivalentes, suites négligeables}
\cpge{mpsi-mp2i.nombres-complexes.b, pcsi.nombres-complexes.b, ptsi.nombres-complexes.b, mpsi-mp2i.nombres-reels-suites.h, pcsi.nombres-reels-suites.f, ptsi.nombres-reels-suites.f, mpsi-mp2i.nombres-reels-suites.i, pcsi.nombres-reels-suites.g, ptsi.nombres-reels-suites.g, bcpst1.nombres-complexes.a, bcpst1.suites-reelles}
\module{Analyse}
\niveau{L1}
\difficulte{}

\contenu{
\question{
Soit $f:\Cc\rightarrow \Cc$ d\'{e}finie par :
$$\forall z\in \Cc,\text{ }f(z)=\frac{z+\left| z\right| }{2}.$$
\'Etudier la suite d\'{e}finie par :
$$z_{0}\in \Cc,\forall n\in \Nn,z_{n+1}=f(z_{n}).$$
\emph{Indication} : on \'{e}crira $z_{n}=\rho _{n}e^{i\phi _{n}}, $o\`{u}
$(\rho_{n},\phi _{n})\in \Rr^{+*}\times ]-\pi ,\pi [$ et on utilisera :
$$\sin \phi =2^{n}\sin \frac{\phi }{2^{n}}\prod\limits_{i=1}^{n}\cos \frac{\phi }{2^{i}}.$$
}
}
